\chapter{Portail back-office de validation interne}

\section{Introduction}
% TODO: Introduire le role interne du portail back-office pour Smollan.

\section{Objectif du portail de validation}
% TODO: Expliquer que ce portail porte les controles internes et le suivi des anomalies.

\section{Moteur de regles de validation}
% TODO: Presenter le moteur Python et les regles enregistrees.

\section{Tables validation\_run\_log et validation\_results}
% TODO: Decrire la tracabilite des executions et des anomalies detectees.

\section{Suivi des anomalies OSA et GPS}
% TODO: Presenter les regles OSA et GPS sans exagerer leur maturite.

\section{Dashboard de validation existant}
\begin{figure}[htbp]
    \centering
    \fbox{\parbox{0.82\textwidth}{\centering Espace reserve a la capture du portail back-office de validation.}}
    \caption{Dashboard interne de validation}
    \label{fig:dashboard_validation}
\end{figure}

\section{Workflow de validation prevu}
% TODO: Presenter comme perspective : pending/reviewed, commentaires, action needed, reviewed_by, reviewed_at.

\section{Conclusion}
% TODO: Faire la transition vers les resultats, limites et perspectives.

